\scanpage{42}\mspage{41}
c) L'ens des fermés pour $\tau'$ est en corr. biuniv. avec l'ens des couples
\[
\begin{aligned}
(B,A_{u}) \qquad & B\subsetT\ \text{fermé pour }\tau'_0\\
& A_u\ \text{fermé de }U=S\setminus T\ \text{pour }\tau\mid(S\setminus T)
\end{aligned}
\]
\[
\text{tels que}\qquad \overline{A_u}\cap T\subset B \qquad
\text{(adh. prise pour }\tau\text{)}\ \deleted{\tau)}.
\]

d) Si $\tau'_0$ et $\tau$ ultra topologiques, $\tau'$ aussi. La relation de préordre associée \deleted{se définit} \inserted{se décrit} ainsi. Pour $x,y\in S$, on a $x\underset{\tau'}{\leq}y$ sss
\[
\left\{\begin{aligned}
&\text{ou bien}\quad x,y\in T,\qquad x\underset{\tau'}{\leq}y\\
&\left.\begin{aligned}
&\text{ou bien}\quad x,y\in U=S\setminus T,\qquad x\underset{\tau}{\leq}y\\
&\text{ou bien}\quad x\in T,\ y\in U,\qquad x\underset{\tau}{\leq}y
\end{aligned}\right]
\end{aligned}\right.
\]
\margin{en d'autres termes, si $y\in U$, on a $x\underset{\tau'}{\leq}y$ sss \unsure{a} $x\underset{\tau}{\leq}y$.}

e) \inserted{Sous les conditions de d)} Considérons les ens ordonnés $I,I'$ associés à $\tau,\tau'$, et la partie \inserted{fermée} $J\subset I$ correspondant à la partie fermée $T$ de $(S,\tau)$, enfin l'application $\varphi:I'\mathbin{\text{---}}I$ canonique.
\diagram{\begin{tikzcd}[ampersand replacement=\&,column sep=small,row sep=large]
\varphi^{-1}(J)={} \& J' \arrow[d] \& \subset \& I' \arrow[d,"\varphi"] \& \overline{\tau'} \arrow[d] \\
{} \& J \& \subset \& I \& \overline{\tau}
\end{tikzcd}}

\deleted{Mais}, \deleted{\unsure{la|le}} Soit $\overline{\tau'}$ \deleted{est} la top de $\tau'$ déduite de $\tau'$ par passage au quotient (ou encore, associée à la relation \uline{d'ordre} sur $I'$ associé au préordre), $\widetilde{\tau}$ la topologie sur $I'$ image inverse de $\overline{\tau}$. Alors $\overline{\tau'}$ est la moins fine des top. sur $I'$ qui \inserted{sont plus fines que $\widetilde{\tau}$ et} induisent sur $J'$ la $=$ \deleted{top} topologie $\overline{\tau}\mid J'$. Elle est
